\chapter{Antecedentes}
\label{cap:antecedentes}

\section{Estado del Arte en Domótica}

La domótica ha experimentado en la última década una rápida evolución, impulsada principalmente por la proliferación de dispositivos conectados de bajo coste. Plataformas como Amazon Alexa, Google Home o Apple HomeKit han democratizado el acceso a la automatización del hogar, pero lo han hecho a costa de crear ecosistemas cerrados y dependientes de la conectividad a Internet y de la voluntad comercial de sus fabricantes.

Como alternativa \textit{open-source}, proyectos como \textbf{Home Assistant}~\cite{home_assistant_2019} han ganado una considerable tracción en la comunidad, demostrando que es viable gestionar múltiples protocolos IoT desde un servidor local. Home Assistant soporta más de 2000 integraciones y se ejecuta en hardware modesto como una Raspberry Pi. Sin embargo, su modelo es generalista y no incorpora de forma nativa lógica de toma de decisiones basada en Inteligencia Artificial. Las automatizaciones se definen mediante reglas YAML, no mediante modelos de lenguaje que entiendan contexto.

\section{Edge AI y Modelos de Lenguaje Locales}

La madurez alcanzada por herramientas como \textbf{Ollama}~\cite{ollama} ha sido el catalizador tecnológico que hace viable la fase actual del proyecto. Ollama permite ejecutar modelos de lenguaje de tamaño medio (7B--14B parámetros) en hardware de consumo con rendimiento aceptable. Esto elimina la dependencia de APIs externas de pago como OpenAI o Google Gemini para la lógica cognitiva del sistema, haciendo que la promesa de soberanía del dato sea completamente realizable.

El panorama de modelos pequeños ha evolucionado rápidamente. La familia \textbf{Qwen2.5}~\cite{hui_qwen25_coder_2024}, de Alibaba Cloud, ha demostrado que modelos con 1.5B y 7B parámetros pueden competir en tareas específicas con modelos mucho mayores. El modelo de 1.5B ocupa menos de 2~GB en RAM y responde en cientos de milisegundos en CPU, mientras que la versión de 7B (que cabe en 6~GB) mantiene una calidad conversacional equiparable a modelos de 30~B de generaciones anteriores.

\subsection{TinyML y Procesamiento en el Edge}

El campo del TinyML~\cite{kallimani_tinyml_2023} ha demostrado que es posible ejecutar modelos de machine learning en microcontroladores con apenas kilobytes de memoria. Aunque este proyecto no llega a ese extremo (los modelos se ejecutan en un servidor x86), los principios del TinyML (minimización de dependencias, optimización para hardware limitado, procesamiento local de datos) informan el diseño del sistema.

\subsection{Asistentes de Voz Locales}

Actualmente es posible construir sistemas de reconocimiento y síntesis de voz que funcionen completamente desconectados de Internet. Aunque la calidad sigue siendo inferior a la de los asistentes comerciales (que pueden aprovechar modelos con cientos de miles de millones de parámetros), la brecha se reduce año tras año. La decisión de este proyecto de no incorporar entrada por voz de momento responde a esta limitación objetiva.

\section{Fase Previa del Proyecto: Prototipado Remoto}

Durante la primera fase de este proyecto, y como punto de partida exploratorio, se utilizó la infraestructura de computación de la Universidad de Córdoba (UCO). El acceso a servidores con GPU RTX mediante VPN y SSH permitió prototipar y validar la arquitectura de agentes OpenClaw sin la necesidad de contar desde el inicio con hardware propio de altas prestaciones, lo que supuso un ahorro económico significativo y aceleró el ciclo de validación de hipótesis.

Esta fase permitió validar el enfoque agéntico y sirvió de base para las decisiones de diseño que definen la arquitectura actual. Las lecciones aprendidas durante el prototipado remoto quedaron registradas como registros de decisión (ADR) y se tuvieron en cuenta al migrar el sistema al entorno completamente local de la segunda fase.

Sobre la base de esta experiencia, se seleccionó el orquestador principal del sistema, así como las herramientas y librerías que lo componen. En paralelo, se evaluó el ecosistema \textit{open-source} de IoT (ESPHome, Tasmota, OpenMQTTGateway) como posible base para las integraciones, pero se optó por un enfoque más ligero basado en \textit{drivers} Python específicos para cada dispositivo, priorizando la simplicidad y el control sobre el hardware.

\section{Orquestadores de Agentes: OpenClaw y Alternativas}

El orquestador principal del ecosistema es \textbf{OpenClaw}~\cite{openclaw}, un \textit{framework} agéntico para Node.js que proporciona canales de usuario (Telegram, web), ejecución de herramientas y coordinación de modelos de lenguaje. Se eligió frente a otras opciones por su ligereza, simplicidad de configuración y modularidad.

Durante la fase de investigación se analizó \textbf{Hermes}~\cite{hermes_agent}, una plataforma emergente con objetivos similares. Hermes incorpora un mecanismo de autoaprendizaje que permite al sistema ajustar su comportamiento basándose en la retroalimentación del usuario sin intervención manual. OpenClaw se mantuvo como opción principal por su madurez y ecosistema, pero la idea de incorporar autoaprendizaje se adoptó como una línea de desarrollo propia, implementada en el Smart Advisor mediante la detección de patrones de uso y ajuste dinámico de zonas de confort.

\section{Decisiones sobre Infraestructura de Datos}

Durante el diseño se evaluó la posibilidad de incorporar una \textbf{base de datos vectorial} para mejorar la recuperación semántica de comandos y contextos. Se consideraron tres opciones: \texttt{hnswlib} (librería ligera sin servidor), \texttt{ChromaDB} (base de datos vectorial embebida) y \texttt{Qdrant} (motor especializado con alta eficiencia).

La decisión fue \textbf{no incorporar ninguna} en esta fase. El razonamiento fue doble: primero, el volumen de datos del ecosistema actual es reducido (15 intenciones clasificadas, menos de 200 entradas en caché LRU), y una base de datos vectorial añadiría una capa de complejidad y consumo de RAM no justificados por la ganancia de rendimiento. Segundo, los \textit{embeddings} de \texttt{bge-m3}~\cite{bge_m3} precalculados al arranque y la caché LRU existente cubren sobradamente las necesidades actuales. Queda como línea futura para cuando el sistema escale a decenas de dispositivos y cientos de comandos.

\section{Resumen del Contexto Tecnológico}

En síntesis, el momento tecnológico actual ofrece las piezas necesarias para construir un sistema de domótica local con IA en el \textit{edge}: modelos de lenguaje pequeños pero capaces (\texttt{qwen2.5}), herramientas de inferencia local maduras (Ollama), \textit{frameworks} agénticos ligeros (OpenClaw), librerías de integración IoT que funcionan sin nube (\texttt{tinytuya}, \texttt{bleak}) y una comunidad activa de ingeniería inversa que permite domesticar dispositivos cerrados. La novedad de este proyecto no está en ninguna de estas piezas por separado, sino en su integración en un sistema cohesivo que mantiene la soberanía del usuario como principio de diseño.
